\section{Giới thiệu (Introduction)}
\label{sec:introduction}

\subsection{Bối cảnh và Tầm quan trọng của Bài toán STLF}
Dự báo phụ tải điện ngắn hạn (\textit{Short-Term Load Forecasting} -- STLF), với khoảng thời gian dự báo từ vài giờ đến vài ngày tới, là bài toán trung tâm trong công tác quản lý và vận hành hệ thống năng lượng hiện đại~\cite{hong2016probabilistic}. Trong một lưới điện xoay chiều, cân bằng công suất tức thời giữa tổng nguồn phát và tổng phụ tải tiêu thụ:
\begin{equation}
    \sum_{i} P_{\text{gen}, i}(t) = \sum_{j} P_{\text{load}, j}(t) + P_{\text{loss}}(t)
\end{equation}
là điều kiện tiên quyết để duy trì tần số danh định (50/60~Hz) và đảm bảo tính ổn định động học của toàn hệ thống. Do năng lượng điện quy mô lớn rất khó và tốn kém để lưu trữ bằng pin tích năng, mọi sai số trong dự báo phụ tải đều kéo theo tổn thất kinh tế khổng lồ: nếu dự báo thiếu, đơn vị vận hành hệ thống điện (\textit{Independent System Operator} -- ISO) buộc phải huy động khẩn cấp các tổ máy phát điện chạy dầu/khí dự phòng (\textit{peaking units}) với giá thị trường giao ngay (\textit{spot price}) cao gấp nhiều lần; ngược lại, nếu dự báo thừa, các nhà máy điện cơ bản buộc phải giảm tải non công suất, gây lãng phí nhiên liệu và phát thải khí nhà kính không đáng có. Nghiên cứu của Hong và Fan~\cite{hong2016probabilistic} đã chỉ ra rằng việc giảm chỉ $1\%$ sai số dự báo phụ tải trong một hệ thống điện cấp tiểu bang có thể tiết kiệm hàng chục triệu USD chi phí vận hành mỗi năm.

Trong bối cảnh hiện nay, quá trình chuyển dịch năng lượng và hiện tượng biến đổi khí hậu toàn cầu đang đẩy bài toán STLF lên một mức độ phức tạp chưa từng có. Một mặt, tỷ trọng thâm nhập cao của các nguồn năng lượng tái tạo phân tán, đặc biệt là điện mặt trời mái nhà (\textit{rooftop solar PV}), làm biến dạng sâu sắc đồ thị phụ tải thực nhận tại các trạm biến áp (tạo nên hiện tượng đường cong con vịt -- \textit{duck curve}). Mặt khác, sự gia tăng tần suất của các hình thái thời tiết cực đoan như các đợt sóng nhiệt kéo dài (\textit{heatwaves}) hoặc các đợt rét đậm làm bùng nổ nhu cầu sử dụng các thiết bị làm mát và sưởi ấm (\textit{Heating, Ventilation, and Air Conditioning} -- HVAC), đẩy phụ tải đỉnh vượt xa các ngưỡng lịch sử thông thường.

\subsection{Cuộc thi IISE PG\&E 2025 và Đóng góp của Bài báo Gốc}
Để thúc đẩy các giải pháp học máy tiên tiến đối phó với thách thức trên, Viện Kỹ sư Công nghiệp và Hệ thống (\textit{Institute of Industrial and Systems Engineers} -- IISE) phối hợp cùng Công ty Điện lực và Khí đốt Thái Bình Dương (\textit{Pacific Gas and Electric Company} -- PG\&E) đã tổ chức cuộc thi \textit{IISE PG\&E Energy Analytics Challenge 2025}~\cite{pge2025zenodo}. Bộ dữ liệu cung cấp chuỗi thời gian phụ tải thực tế theo từng giờ trên địa bàn quản lý của PG\&E thuộc bang California trong vòng 3 năm liên tục (2020--2022), kèm theo các thông số khí tượng đo đạc đồng bộ tại 5 trạm quan trắc địa lý trọng yếu, bao gồm nhiệt độ mặt đất và bức xạ mặt trời trực tiếp (\textit{Global Horizontal Irradiance} -- GHI).

Tại cuộc thi này, công trình nghiên cứu của Roy, Pyltsov và Hu~\cite{roy2025hourly} (arXiv:2505.11390) đã tạo nên một dấu ấn học thuật sâu sắc. Trong khi xu hướng nghiên cứu thịnh hành đang dồn sự chú ý vào các mô hình \textit{Deep Learning} phức tạp với hàng triệu tham số dựa trên kiến trúc \textit{Transformer} (như PatchTST~\cite{nie2023time}, Informer~\cite{zhou2021informer}, Autoformer~\cite{wu2021autoformer}), nhóm tác giả đã chứng minh một luận điểm thực nghiệm mạnh mẽ: \textbf{các mô hình hồi quy dạng bảng phân đoạn theo giờ (24 mô hình XGBoost độc lập cho 24 giờ trong ngày) hoàn toàn đánh bại các kiến trúc Transformers phức tạp}. Cụ thể, bài báo gốc đạt sai số RMSE $170.84$~MW và MAPE $5.38\%$ trên tập kiểm tra độc lập năm thứ ba (Year 3), vượt trội hơn toàn bộ các mô hình học sâu chuỗi thời gian lớn tham gia benchmark. Bí quyết cốt lõi của công trình nằm ở việc:
\begin{enumerate}
    \item Ứng dụng kỹ thuật phân tích thành phần chính (\textit{Principal Component Analysis} -- PCA) trên 5 trạm nhiệt độ và 5 trạm GHI để nén thông tin và triệt tiêu triệt để hiện tượng đa cộng tuyến (\textit{multicollinearity}).
    \item Phân đoạn không gian bài toán thành 24 mô hình XGBoost riêng biệt, giúp giải phóng mô hình khỏi gánh nặng phải học sự biến thiên nhật triều (\textit{diurnal variation}) vốn rất phức tạp giữa các khung giờ ngày và đêm.
    \item Bổ sung các đặc trưng biến trễ thời gian 1 giờ của thành phần chính (\textit{PCA Lag1}).
\end{enumerate}

\subsection{Khoảng trống Nghiên cứu (Research Gaps)}
Mặc dù bài báo của Roy và cộng sự~\cite{roy2025hourly} thiết lập một tiêu chuẩn thực nghiệm rất cao, quá trình phân tích và tái lập sâu sắc của chúng tôi đã chỉ ra ba khoảng trống phương pháp luận trọng yếu:

\begin{enumerate}
    \item \textbf{Thiếu vắng tri thức nhiệt động lực học công trình (\textit{Absence of Thermodynamic Inductive Bias}):} Phụ tải điện đô thị chịu sự chi phối nặng nề nhất bởi hành vi bật/tắt thiết bị nhiệt HVAC của con người. Mối quan hệ giữa nhiệt độ môi trường và công suất điện không bao giờ là một hàm tuyến tính đơn điệu, mà là một đường cong chữ U (\textit{U-shaped curve}) với ngưỡng kích hoạt vật lý rõ ràng: khi nhiệt độ thấp hơn khoảng $18^\circ\text{C}$, nhu cầu sưởi ấm bắt đầu tăng; khi nhiệt độ vượt quá $18^\circ\text{C}$, nhu cầu điều hòa không khí bùng nổ phi tuyến tính. Việc chỉ dựa vào các thành phần chính PCA tuyến tính thuần túy đã vô tình làm phẳng các ngưỡng phi tuyến nhạy cảm này, khiến mô hình gặp khó khăn trong việc nắm bắt các phụ tải đỉnh cực đoan vào các đợt nắng nóng mùa hè.
    \item \textbf{Sự ngắt quãng giả tạo của các đặc trưng lịch rời rạc (\textit{Calendar Discontinuity}):} Bài báo gốc sử dụng 12 biến giả nhị phân rời rạc (\textit{one-hot dummy variables}) để đại diện cho 12 tháng trong năm. Cách tiếp cận này tạo ra các bước nhảy bậc thang gián đoạn tại thời điểm giao mùa (ví dụ giữa ngày 31 tháng 12 và ngày 01 tháng 01), hoàn toàn mâu thuẫn với quy luật thiên văn học của Trái Đất nơi vị trí mặt trời và góc chiếu bức xạ biến thiên hoàn toàn trơn và tuần hoàn liên tục.
    \item \textbf{Hạn chế của mô hình đơn nhất và sự bỏ ngỏ đa dạng sai số (\textit{Lack of Regime-Aware Model Synergy}):} Bài báo gốc chỉ xem xét các họ mô hình một cách biệt lập (chỉ dùng thuần XGBoost hoặc thuần Transformer). Trên thực tế, các kiến trúc mô hình khác nhau sở hữu thiên kiến quy nạp (\textit{inductive bias}) hoàn toàn khác nhau. Mô hình tuyến tính có chính quy hóa (như ElasticNet) có tính khái quát hóa cực kỳ ổn định trên các vùng phụ tải nền đều đặn, ít biến động; ngược lại, các mô hình cây quyết định phân nhánh (như XGBoost) lại cực kỳ nhạy bén trong việc bắt trọn các quan hệ phi tuyến phức tạp trong những giờ cao điểm. Việc không kết hợp tính bù trừ sai số giữa các họ mô hình này theo từng chế độ vận hành lưới điện là một sự lãng phí tài nguyên mô hình hóa to lớn.
\end{enumerate}

\subsection{Mục tiêu và Đóng góp Cốt lõi của Nhóm Nghiên cứu}
Nhằm khắc phục toàn diện các hạn chế trên, nghiên cứu này mở rộng bài toán STLF của PG\&E thành một khung kiến trúc ensemble đa tầng nhận biết ngữ cảnh hoàn chỉnh. Cụ thể, các đóng góp khoa học nổi bật của nghiên cứu bao gồm:

\begin{itemize}
    \item \textbf{Tái lập độc lập và chuẩn hóa thực nghiệm:} Chúng tôi tái lập thành công và kiểm chứng toàn diện toàn bộ quy trình thí nghiệm, cấu hình tham số và kết quả của bài báo gốc (Roy et al., 2025) trên cả ba giao thức kiểm chứng chéo và tập test năm thứ ba, thiết lập một nền tảng đối chuẩn vững chắc.
    \item \textbf{Đề xuất kỹ nghệ đặc trưng định hướng vật lý (\textit{Physics-Informed Feature Engineering}):} Nhóm tích hợp chỉ số \textit{Heating Degree Days} (HDD) và \textit{Cooling Degree Days} (CDD) dựa trên ngưỡng kích hoạt vật lý $18^\circ\text{C}$ theo tiêu chuẩn ASHRAE~\cite{ashrae2021fundamentals}, kết hợp cùng chuỗi điều hòa \textit{Fourier harmonics} ($\sin, \cos$) liên tục theo ngày và theo năm. Kết quả cho thấy chỉ riêng việc bổ sung các đặc trưng này vào mô hình 24-XGBoost đã giúp giảm ngay sai số RMSE từ $171.66$~MW xuống $160.65$~MW (giảm $11.01$~MW so với baseline tái lập, và giảm $10.19$~MW so với baseline công bố $170.84$~MW của bài báo gốc) trên tập kiểm tra độc lập Year 3.
    \item \textbf{Khám phá và lượng hóa tính bù trừ sai số đa chế độ (\textit{Dual-Regime Discovery}):} Qua phân tích tương quan phần dư (\textit{residual correlation}), chúng tôi chứng minh hệ số tương quan sai số giữa ElasticNet và XGBoost chỉ ở mức $r = 0.6131$. Đặc biệt, chúng tôi phát hiện sự phân tầng rõ rệt: ElasticNet vượt trội hơn XGBoost trong chế độ phụ tải nền ban đêm (\textit{nocturnal baseload}, 00:00--06:00: RMSE $131.33$~MW so với $141.37$~MW), trong khi XGBoost vượt trội trong chế độ quá nhiệt mùa hè (\textit{summer heatwaves}, CDD $> 3^\circ\text{C}$: RMSE $244.16$~MW so với $313.27$~MW).
    \item \textbf{Đề xuất kiến trúc Context-Aware Stacking (CAS):} Nhóm phát triển kiến trúc Stacking hai tầng với mô hình \textit{meta-learner} LightGBM được cung cấp đồng thời cả dự báo của các mô hình cơ sở (\textit{base models}) và vector ngữ cảnh khí tượng -- thời gian thực (\textit{real-time environmental context}). CAS hoạt động như một bộ định tuyến thích nghi động (\textit{dynamic router}), đạt kỷ lục mới với RMSE \textbf{141.81~MW}, MAPE \textbf{4.75\%}, và $R^2$ \textbf{0.9136} trên tập kiểm tra Year 3 (giảm \textbf{$-$17.0\% RMSE} so với công bố của bài báo gốc).
    \item \textbf{Khả năng giải thích và phân tích theo giờ qua SHAP:} Chúng tôi cung cấp biểu đồ nhiệt SHAP 24 giờ trực quan hóa sự dịch chuyển động học của tầm quan trọng đặc trưng (\textit{feature importance}) trong suốt chu kỳ ngày đêm, cung cấp cơ sở vững chắc cho các kỹ sư vận hành lưới điện.
\end{itemize}
